\documentclass[aspectratio=169]{beamer}

\usetheme{Madrid}
\usecolortheme{default}

\usepackage[UTF8]{ctex}
\usepackage{amsmath,amssymb}
\usepackage{booktabs}
\usepackage{hyperref}
\usepackage{listings}
\usepackage{xcolor}

\definecolor{codebg}{RGB}{245,245,245}
\definecolor{codekw}{RGB}{0,0,150}
\definecolor{codestr}{RGB}{163,21,21}
\definecolor{codecom}{RGB}{0,128,0}

\lstset{
  backgroundcolor=\color{codebg},
  basicstyle=\ttfamily\small,
  keywordstyle=\color{codekw}\bfseries,
  stringstyle=\color{codestr},
  commentstyle=\color{codecom},
  showstringspaces=false,
  frame=single,
  breaklines=true,
  columns=fullflexible
}

\title{Shell 编程基础}
\author{《人工智能与数学软件》课程}
\date{\today}

\begin{document}

\begin{frame}
  \titlepage
\end{frame}

\begin{frame}{本节目标}
学完这一节后，希望能够：
\begin{itemize}
  \item 理解什么是 Shell 脚本，为什么要写脚本；
  \item 会写最基本的 \texttt{bash} 脚本；
  \item 掌握变量、参数、分支、循环的基本用法；
  \item 能把若干条命令组织成一个可复用的小程序；
  \item 能读懂并修改简单的系统管理脚本。
\end{itemize}
\end{frame}

\section{为什么要写 Shell 脚本}

\begin{frame}{为什么要写 Shell 脚本}
Shell 编程的核心思想是：

\begin{center}
\Large
把“手工重复输入的命令”写成“可以反复执行的程序”
\end{center}

\vspace{0.5em}
常见用途：
\begin{itemize}
  \item 批量处理文件；
  \item 自动备份、自动清理；
  \item 自动检查系统状态；
  \item 把多个命令组织成一个工作流。
\end{itemize}
\end{frame}

\begin{frame}{Shell 编程的特点}
优点：
\begin{itemize}
  \item 直接调用系统命令；
  \item 写起来快，适合自动化；
  \item 非常适合文件、进程、日志、网络等任务。
\end{itemize}

\vspace{0.5em}
局限：
\begin{itemize}
  \item 大型程序维护性较差；
  \item 数据结构和语法能力有限；
  \item 复杂算法通常更适合 Python、C/C++ 等语言。
\end{itemize}

\vspace{0.5em}
所以，Shell 最适合做“\alert{系统胶水层}”。
\end{frame}

\section{第一个脚本}

\begin{frame}[fragile]{第一个 Shell 脚本}
先写一个最简单的脚本文件 \texttt{hello.sh}：

\begin{lstlisting}[language=bash]
#!/bin/bash
echo "Hello, Shell"
echo "Today is:"
date
\end{lstlisting}

说明：
\begin{itemize}
  \item 第一行 \texttt{\#!/bin/bash} 称为 shebang；
  \item 它告诉系统：用 \texttt{bash} 来解释这个脚本；
  \item 后面每一行本质上都是普通 Shell 命令。
\end{itemize}
\end{frame}

\begin{frame}[fragile]{如何运行脚本}
方法一：直接交给 bash 执行
\begin{lstlisting}[language=bash]
bash hello.sh
\end{lstlisting}

方法二：先加执行权限，再直接运行
\begin{lstlisting}[language=bash]
chmod u+x hello.sh
./hello.sh
\end{lstlisting}

教学上应强调：
\begin{itemize}
  \item 脚本本质上是文本文件；
  \item “能不能执行”取决于权限和解释器；
  \item \texttt{./} 表示“当前目录中的脚本”。
\end{itemize}
\end{frame}

\begin{frame}[fragile]{脚本中的注释}
注释以 \texttt{\#} 开头：

\begin{lstlisting}[language=bash]
#!/bin/bash
# 这是一个简单示例
echo "Start"
# 输出当前用户
whoami
\end{lstlisting}

作用：
\begin{itemize}
  \item 解释脚本功能；
  \item 标注关键步骤；
  \item 便于自己和别人维护。
\end{itemize}
\end{frame}

\section{变量与参数}

\begin{frame}[fragile]{变量的基本写法}
Shell 中定义变量时，等号两边不要有空格：

\begin{lstlisting}[language=bash]
name="Alice"
course="Shell Programming"
echo $name
echo $course
\end{lstlisting}

说明：
\begin{itemize}
  \item 使用变量时，在前面加 \texttt{\$}；
  \item 右边通常写成字符串；
  \item 变量名一般用字母、数字、下划线。
\end{itemize}
\end{frame}

\begin{frame}[fragile]{一个变量示例}
\begin{lstlisting}[language=bash]
#!/bin/bash
user=$(whoami)
today=$(date +%F)

echo "Current user: $user"
echo "Today: $today"
\end{lstlisting}

这里用到了命令替换：
\begin{itemize}
  \item \texttt{\$(whoami)}：把命令输出赋给变量；
  \item \texttt{\$(date +\%F)}：得到类似 2026-03-30 的日期。
\end{itemize}
\end{frame}

\begin{frame}[fragile]{位置参数：让脚本接收输入}
脚本运行时可以带参数：

\begin{lstlisting}[language=bash]
bash greet.sh Tom
\end{lstlisting}

脚本内容：
\begin{lstlisting}[language=bash]
#!/bin/bash
echo "Hello, $1"
\end{lstlisting}

常见位置参数：
\begin{itemize}
  \item \texttt{\$1}, \texttt{\$2}, \texttt{\$3}：第 1、2、3 个参数；
  \item \texttt{\$0}：脚本名；
  \item \texttt{\$\#}：参数个数；
  \item \texttt{\$@}：所有参数。
\end{itemize}
\end{frame}

\begin{frame}[fragile]{参数示例}
\begin{lstlisting}[language=bash]
#!/bin/bash
echo "Script name: $0"
echo "First arg: $1"
echo "Second arg: $2"
echo "Number of args: $#"
echo "All args: $@"
\end{lstlisting}

例如执行：
\begin{lstlisting}[language=bash]
bash args.sh file1.txt file2.txt
\end{lstlisting}
\end{frame}

\section{条件判断}

\begin{frame}[fragile]{if 条件语句}
Shell 中最常见的分支结构是 \texttt{if}：

\begin{lstlisting}[language=bash]
if [ 条件 ]
then
  命令
else
  命令
fi
\end{lstlisting}

注意：
\begin{itemize}
  \item \texttt{[ ]} 两边要留空格；
  \item \texttt{fi} 是 \texttt{if} 的结束；
  \item 条件成立时执行 \texttt{then} 分支。
\end{itemize}
\end{frame}

\begin{frame}[fragile]{判断文件是否存在}
\begin{lstlisting}[language=bash]
#!/bin/bash
if [ -f "$1" ]
then
  echo "File exists."
else
  echo "File does not exist."
fi
\end{lstlisting}

执行：
\begin{lstlisting}[language=bash]
bash check_file.sh notes.txt
\end{lstlisting}

这里：
\begin{itemize}
  \item \texttt{-f} 判断是否为普通文件；
  \item \texttt{"\$1"} 表示脚本的第一个参数。
\end{itemize}
\end{frame}

\begin{frame}[fragile]{常见测试条件}
\begin{center}
\begin{tabular}{ll}
\toprule
写法 & 含义 \\
\midrule
\texttt{-f file} & 是否为普通文件 \\
\texttt{-d dir} & 是否为目录 \\
\texttt{-e path} & 路径是否存在 \\
\texttt{-z str} & 字符串是否为空 \\
\texttt{a = b} & 字符串是否相等 \\
\texttt{n -gt m} & 是否大于 \\
\texttt{n -lt m} & 是否小于 \\
\bottomrule
\end{tabular}
\end{center}

例如：
\begin{lstlisting}[language=bash]
if [ "$1" = "yes" ]
then
  echo "confirmed"
fi
\end{lstlisting}
\end{frame}

\section{循环}

\begin{frame}[fragile]{for 循环}
\texttt{for} 循环适合遍历一组对象：

\begin{lstlisting}[language=bash]
for x in a b c
do
  echo $x
done
\end{lstlisting}

输出：
\begin{lstlisting}[language=bash]
a
b
c
\end{lstlisting}
\end{frame}

\begin{frame}[fragile]{for 循环处理多个文件}
\begin{lstlisting}[language=bash]
#!/bin/bash
for f in *.txt
do
  echo "Found file: $f"
done
\end{lstlisting}

用途：
\begin{itemize}
  \item 批量处理同类文件；
  \item 是 Shell 自动化中最常见的模式之一。
\end{itemize}
\end{frame}

\begin{frame}[fragile]{while 循环}
\texttt{while} 循环适合“条件满足就继续执行”：

\begin{lstlisting}[language=bash]
#!/bin/bash
count=1
while [ $count -le 3 ]
do
  echo $count
  count=$((count + 1))
done
\end{lstlisting}

说明：
\begin{itemize}
  \item \texttt{\$((...))} 用于整数运算；
  \item 这里会输出 1、2、3。
\end{itemize}
\end{frame}

\section{实用例子}

\begin{frame}[fragile]{例 1：批量统计文本文件行数}
脚本：
\begin{lstlisting}[language=bash]
#!/bin/bash
for f in *.txt
do
  echo -n "$f : "
  wc -l "$f"
done
\end{lstlisting}

这个脚本做的事：
\begin{itemize}
  \item 找到当前目录下所有 \texttt{.txt} 文件；
  \item 逐个调用 \texttt{wc -l}；
  \item 输出每个文件的行数。
\end{itemize}
\end{frame}

\begin{frame}[fragile]{例 2：简单系统体检脚本}
\begin{lstlisting}[language=bash]
#!/bin/bash
echo "User: $(whoami)"
echo "Kernel:"
uname -a
echo
echo "Memory:"
free -h
echo
echo "Disk:"
df -h
\end{lstlisting}

用途：
\begin{itemize}
  \item 快速查看当前用户；
  \item 查看内核信息；
  \item 查看内存和磁盘状态。
\end{itemize}
\end{frame}

\begin{frame}[fragile]{例 3：检查参数并给出提示}
\begin{lstlisting}[language=bash]
#!/bin/bash
if [ $# -lt 1 ]
then
  echo "Usage: bash greet.sh NAME"
  exit 1
fi

echo "Hello, $1"
\end{lstlisting}

这里增加了一个重要观念：
\begin{itemize}
  \item 脚本应该检查输入是否正确；
  \item \texttt{exit 1} 表示异常退出。
\end{itemize}
\end{frame}

\section{要点总结}

\begin{frame}{Shell 编程的几个关键要点}
\begin{enumerate}
  \item 脚本就是一串按顺序执行的 Shell 命令；
  \item shebang 指定解释器；
  \item 变量和参数让脚本更灵活；
  \item \texttt{if} 负责分支，\texttt{for}/\texttt{while} 负责循环；
  \item Shell 最适合自动化系统任务和批量处理。
\end{enumerate}
\end{frame}

\begin{frame}{这一节最值得记住的模式}
建议记住以下 4 个“模板”：

\vspace{0.5em}
\textbf{1. 脚本头}
\begin{itemize}
  \item \texttt{\#!/bin/bash}
\end{itemize}

\textbf{2. 变量}
\begin{itemize}
  \item \texttt{name="Alice"}，使用时 \texttt{\$name}
\end{itemize}

\textbf{3. 判断}
\begin{itemize}
  \item \texttt{if [ -f "\$1" ] ... fi}
\end{itemize}

\textbf{4. 遍历}
\begin{itemize}
  \item \texttt{for f in *.txt; do ...; done}
\end{itemize}
\end{frame}

\section{正则表达式与 Shell}

\begin{frame}{正则表达式与 Shell}
\centering
\Large
用模式匹配处理文本
\end{frame}

\begin{frame}{为什么要学正则表达式}
Shell 编程中，经常需要处理文本：

\begin{itemize}
  \item 从日志中找某些行
  \item 从文件中筛选包含特定模式的内容
  \item 批量提取用户名、日期、数字等信息
\end{itemize}

\vspace{0.8em}
这时就需要一种“按模式匹配字符串”的方法：

\begin{center}
\Large
正则表达式（regular expression）
\end{center}
\end{frame}

\begin{frame}{什么是正则表达式}
正则表达式可以理解为：

\begin{center}
\Large
描述一类字符串的规则
\end{center}

\vspace{0.8em}
例如：
\begin{itemize}
  \item \texttt{abc}：匹配字符串 \texttt{abc}
  \item \texttt{[0-9]}：匹配一个数字
  \item \texttt{[a-z]+}：匹配一个或多个小写字母
\end{itemize}

\vspace{0.8em}
在 Shell 中，最常见的使用场景是配合：
\begin{itemize}
  \item \texttt{grep}
  \item \texttt{sed}
  \item \texttt{awk}
\end{itemize}
\end{frame}

\begin{frame}{常见正则表达式符号}
\begin{center}
\begin{tabular}{ll}
\toprule
符号 & 含义 \\
\midrule
\texttt{.} & 任意单个字符 \\
\texttt{[abc]} & 匹配 \texttt{a,b,c} 中任一个 \\
\texttt{[0-9]} & 匹配一个数字 \\
\texttt{[a-z]} & 匹配一个小写字母 \\
\texttt{\^} & 匹配行首 \\
\texttt{\$} & 匹配行尾 \\
\texttt{*} & 前一模式重复 0 次或多次 \\
\texttt{+} & 前一模式重复 1 次或多次 \\
\texttt{?} & 前一模式重复 0 次或 1 次 \\
\bottomrule
\end{tabular}
\end{center}

\vspace{0.5em}
注意：\texttt{+}、\texttt{?} 常见于扩展正则表达式。
\end{frame}

\begin{frame}{读懂几个基本模式}
\begin{itemize}
  \item \texttt{abc}：匹配包含 \texttt{abc} 的文本
  \item \texttt{\^root}：匹配以 \texttt{root} 开头的行
  \item \texttt{bash\$}：匹配以 \texttt{bash} 结尾的行
  \item \texttt{[0-9]}：匹配任意一个数字
  \item \texttt{[0-9][0-9]}：匹配两个数字
  \item \texttt{[a-z]+}：匹配一个或多个连续小写字母
\end{itemize}

\vspace{0.8em}
一个重要观念：
\begin{itemize}
  \item 正则表达式描述的是“模式”，不是某一个固定字符串。
\end{itemize}
\end{frame}

\begin{frame}[fragile]{在 Shell 中最常见的搭配：grep}
\texttt{grep} 用来在文本中查找匹配某个模式的行。

\vspace{0.5em}
基本形式：
\begin{lstlisting}[language=bash]
grep '模式' 文件名
\end{lstlisting}

例如：
\begin{lstlisting}[language=bash]
grep 'root' /etc/passwd
grep '^root' /etc/passwd
grep 'bash$' /etc/passwd
\end{lstlisting}

说明：
\begin{itemize}
  \item 第一条：找包含 \texttt{root} 的行
  \item 第二条：找以 \texttt{root} 开头的行
  \item 第三条：找以 \texttt{bash} 结尾的行
\end{itemize}
\end{frame}

\begin{frame}[fragile]{扩展正则表达式：grep -E}
有些更方便的写法，需要使用扩展正则表达式：

\begin{lstlisting}[language=bash]
grep -E '[0-9]+'
grep -E '^[a-z]+$'
\end{lstlisting}

\vspace{0.5em}
解释：
\begin{itemize}
  \item \verb|[0-9]+|：一个或多个数字
  \item \verb|^[a-z]+$|：整行只包含一个或多个小写字母
\end{itemize}

\vspace{0.8em}
教学上可以简单记成：
\begin{itemize}
  \item \texttt{grep}：基本正则
  \item \texttt{grep -E}：扩展正则，更方便
\end{itemize}
\end{frame}

\begin{frame}[fragile]{例 1：从 \texttt{/etc/passwd} 中筛选}
\textbf{找出以 \texttt{root} 开头的行：}
\begin{lstlisting}[language=bash]
grep '^root' /etc/passwd
\end{lstlisting}

\textbf{找出登录 shell 是 \texttt{/bin/bash} 的用户：}
\begin{lstlisting}[language=bash]
grep 'bash$' /etc/passwd
\end{lstlisting}

\textbf{找出用户名中包含数字的行：}
\begin{lstlisting}[language=bash]
grep -E '[0-9]' /etc/passwd
\end{lstlisting}

\vspace{0.5em}
这说明：正则表达式特别适合处理系统文本文件。
\end{frame}

\begin{frame}[fragile]{例 2：在脚本中筛选文本文件}
下面这个脚本只处理 \texttt{.txt} 文件：

\begin{lstlisting}[language=bash]
#!/bin/bash
for f in *
do
  if echo "$f" | grep -E '\.txt$' > /dev/null
  then
    echo "Text file: $f"
  fi
done
\end{lstlisting}

说明：
\begin{itemize}
  \item \texttt{\.txt\$} 表示以 \texttt{.txt} 结尾
  \item \texttt{> /dev/null} 表示不显示匹配结果，只关心是否匹配成功
\end{itemize}
\end{frame}

\begin{frame}[fragile]{例 3：检查输入是否为整数}
\begin{lstlisting}[language=bash]
#!/bin/bash
if echo "$1" | grep -E '^[0-9]+$' > /dev/null
then
  echo "Input is an integer."
else
  echo "Input is not an integer."
fi
\end{lstlisting}

这个模式：
\begin{lstlisting}[language=bash]
^[0-9]+$
\end{lstlisting}

表示：
\begin{itemize}
  \item 从行首开始
  \item 中间是一个或多个数字
  \item 到行尾结束
\end{itemize}
\end{frame}

\begin{frame}[fragile]{例 4：从文本中找类似 IP 地址的模式}
例如，从日志中找形如 \texttt{a.b.c.d} 的片段：

\begin{lstlisting}[language=bash]
grep -E '[0-9]+\.[0-9]+\.[0-9]+\.[0-9]+' logfile.txt
\end{lstlisting}

说明：
\begin{itemize}
  \item \texttt{\.} 表示字面上的点号
  \item \texttt{[0-9]+} 表示一个或多个数字
\end{itemize}

\vspace{0.5em}
提醒：
\begin{itemize}
  \item 这个模式能匹配“像 IP 的字符串”
  \item 但不保证一定是合法 IP
\end{itemize}
\end{frame}

\begin{frame}[fragile]{grep 与 if 的配合}
在 Shell 中，经常不是要“打印匹配结果”，而是要“根据是否匹配来决定下一步”。

\begin{lstlisting}[language=bash]
if echo "$1" | grep -E '^[a-z]+$' > /dev/null
then
  echo "all lowercase"
else
  echo "not all lowercase"
fi
\end{lstlisting}

这里的逻辑是：
\begin{itemize}
  \item 若 \texttt{grep} 匹配成功，则返回成功状态；
  \item \texttt{if} 就进入 \texttt{then} 分支；
  \item 否则进入 \texttt{else} 分支。
\end{itemize}
\end{frame}

\section{环境变量}

\begin{frame}{什么是环境变量}
环境变量可以理解为：

\begin{center}
\Large
由 Shell 和系统维护、可被程序继承的变量
\end{center}

\vspace{0.8em}
它们常用于保存：
\begin{itemize}
  \item 当前用户信息
  \item 家目录位置
  \item 默认 Shell
  \item 可执行程序搜索路径
  \item 语言、编辑器等配置
\end{itemize}

\vspace{0.8em}
与普通变量相比，环境变量更强调“\alert{会被子进程继承}”。
\end{frame}

\begin{frame}{环境变量举例}
环境变量决定了很多程序运行时的行为。

\vspace{0.5em}
例如：
\begin{itemize}
  \item \texttt{HOME} 决定用户家目录
  \item \texttt{PATH} 决定命令到哪里去找
  \item \texttt{SHELL} 表示当前默认 Shell
  \item \texttt{PWD} 表示当前工作目录
\end{itemize}

\vspace{0.8em}
所以，理解环境变量，就是理解：
\begin{center}
\Large
程序运行时“默认知道些什么”
\end{center}
\end{frame}

\begin{frame}[fragile]{查看环境变量}
查看单个环境变量：

\begin{lstlisting}[language=bash]
echo $HOME
echo $PATH
echo $SHELL
echo $PWD
\end{lstlisting}

查看全部环境变量：

\begin{lstlisting}[language=bash]
env
\end{lstlisting}

或：

\begin{lstlisting}[language=bash]
printenv
\end{lstlisting}
\end{frame}

\begin{frame}{常见环境变量}
\begin{center}
\begin{tabular}{ll}
\toprule
变量名 & 含义 \\
\midrule
\texttt{HOME} & 当前用户的家目录 \\
\texttt{PATH} & 命令搜索路径 \\
\texttt{SHELL} & 当前默认 Shell \\
\texttt{PWD} & 当前工作目录 \\
\texttt{USER} & 当前用户名 \\
\texttt{LANG} & 语言与区域设置 \\
\texttt{EDITOR} & 默认编辑器 \\
\bottomrule
\end{tabular}
\end{center}
\end{frame}

\begin{frame}{特别重要：\texttt{PATH}}
当你输入一条命令，如：

\begin{center}
\texttt{ls}
\end{center}

Shell 会到 \texttt{PATH} 指定的目录中查找对应程序。

\vspace{0.8em}
例如：
\begin{itemize}
  \item \texttt{/usr/bin}
  \item \texttt{/bin}
  \item \texttt{/usr/local/bin}
\end{itemize}

\vspace{0.8em}
这就是为什么直接输入 \texttt{ls} 就能运行，而不必写完整路径。
\end{frame}

\begin{frame}[fragile]{查看 \texttt{PATH} 与命令位置}
\begin{lstlisting}[language=bash]
echo $PATH
which ls
which bash
\end{lstlisting}

\vspace{0.5em}
说明：
\begin{itemize}
  \item \texttt{echo \$PATH}：显示当前搜索路径
  \item \texttt{which ls}：查看 \texttt{ls} 命令实际在哪
  \item \texttt{which bash}：查看 \texttt{bash} 程序路径
\end{itemize}
\end{frame}

\begin{frame}[fragile]{普通变量与环境变量}
定义普通变量：

\begin{lstlisting}[language=bash]
name="Alice"
echo $name
\end{lstlisting}

定义环境变量：

\begin{lstlisting}[language=bash]
export name="Alice"
echo $name
\end{lstlisting}

\vspace{0.8em}
区别：
\begin{itemize}
  \item 普通变量只在当前 Shell 中有效
  \item 环境变量会传递给该 Shell 启动的子进程
\end{itemize}
\end{frame}

\begin{frame}[fragile]{在脚本中使用环境变量}
例如写一个脚本：

\begin{lstlisting}[language=bash]
#!/bin/bash
echo "User: $USER"
echo "Home: $HOME"
echo "Shell: $SHELL"
echo "Current directory: $PWD"
\end{lstlisting}

\vspace{0.8em}
这说明：
\begin{itemize}
  \item 脚本中可以直接使用现成的环境变量；
  \item 它们为脚本提供运行时上下文。
\end{itemize}
\end{frame}

\begin{frame}{环境变量在哪里配置}
常见做法是写入 Shell 的启动文件，例如：

\begin{itemize}
  \item \texttt{\~/.bashrc}
  \item \texttt{\~/.profile}
\end{itemize}

\vspace{0.8em}
例如在 \texttt{\~/.bashrc} 中加入：

\begin{center}
\texttt{export PATH="\$PATH:\$HOME/bin"}
\end{center}

含义是：
\begin{itemize}
  \item 把 \texttt{\$HOME/bin} 加入命令搜索路径
\end{itemize}
\end{frame}

\begin{frame}[fragile]{例子：把自己的脚本目录加入 \texttt{PATH}}
假设你把脚本放在：

\begin{center}
\texttt{\$HOME/bin}
\end{center}

那么在 \texttt{\~/.bashrc} 中加入：

\begin{lstlisting}[language=bash]
export PATH="$PATH:$HOME/bin"
\end{lstlisting}

之后重新加载：

\begin{lstlisting}[language=bash]
source ~/.bashrc
\end{lstlisting}

这样就可以直接运行自己的脚本，而不必每次写 \texttt{./脚本名}。
\end{frame}

\end{document}